\documentclass[10pt,a4paper]{amsart}
\usepackage[margin=19mm]{geometry}
\usepackage{amsmath,amssymb,mathtools}
\usepackage[scr=rsfs,cal=euler]{mathalfa}
\usepackage{xfrac}
\usepackage[unicode,hidelinks,bookmarks=false]{hyperref}
\newcommand{\ie}{i.e.\ }
\newcommand{\cf}{cf.\ }
\newcommand{\resp}{respectively\ }
\newcommand{\Subsection}{\subsection}
\providecommand{\nobreakdash}{\nobreak\hbox{-}}
\setlength{\emergencystretch}{2em}
\multlinegap0pt
\usepackage{bm}
\usepackage{bbm}
\usepackage{dsfont}
\usepackage{xspace}
\usepackage{cancel}
\usepackage{mathtools}
\usepackage{amsthm}
\makeatletter
\@ifundefined{thm}{\newtheorem{thm}{Thm}}{}
\@ifundefined{defn}{\newtheorem{defn}[thm]{Definition}}{}
\@ifundefined{lem}{\newtheorem{lem}[thm]{Lemma}}{}
\@ifundefined{rem}{\newtheorem{rem}[thm]{Remark}}{}
\makeatother
\DeclareMathOperator{\Emp}{\mathcal{E}}%{\mathfrak{m}} %\empirical measure
\DeclareMathOperator{\MT}{\mathcal{M}_{\mathit T} }
\newcommand{\eps}{\varepsilon}
\title[Borel complexity of sets of points with prescribed Birkhoff ]{Borel complexity of sets of points with prescribed Birkhoff averages in Polish dynamical systems with a specification property}
\author{Konrad Deka, Steve Jackson, Dominik Kwietniak, Bill Mance}
\date{Source: arXiv:2210.08937v1 (CC BY 4.0)}
\begin{document}
\maketitle

\noindent\emph{Source and licence.} Borel complexity of sets of points with prescribed Birkhoff averages in Polish dynamical systems with a specification property. Version arXiv:2210.08937v1 (CC BY 4.0), distributed under the Creative Commons Attribution 4.0 International licence, as recorded in the arXiv licence metadata for this exact version and shown on the abstract page https://arxiv.org/abs/2210.08937v1. Full rights notice: licenses/arxiv-2210.08937-CC-BY-4.0.txt.\par\smallskip

\noindent\textit{Editorial scope.} Counted unit: the Lem whose source label is \texttt{lem:main} in the article (source0.tex lines 1421--1432, source label \texttt{lem:main}), delivered together with the article's own complete original proof of it (source0.tex lines 1434--1533, one \texttt{proof} environment of 100 lines). No mathematics is written, repaired or re-derived: statement and proof are the article's own text line for line. Required context retained uncounted: def:faithful2 (lines 697--708); def:saps (lines 711--717); rem:N (lines 720--722); lem:close-segments (lines 1178--1192); lem:protogen (lines 1251--1257). Every definition, display and label of those lines is retained unchanged. Omitted from this package and not redistributed: 1--696, 709--710, 718--719, 723--1177, 1193--1250, 1258--1420, 1433--1433, 1534--2368. Nothing inside a retained excerpt is omitted or altered: every retained body is delivered exactly as the article's lines are written, so no omission receipt of any kind accompanies this package. Citations: the retained excerpts carry no citation command at all, so no bibliography entry of the article is transcribed and no reference is redistributed. Reference adaptation: none was needed and none was made; every reference inside the delivered unit resolves inside it, so no unresolved reference or citation is produced. Printed slips kept literal: no printed word, symbol, hypothesis or number of the counted statement or of its proof was altered. Layout and class: the article's own class is \texttt{amsart} with packages including amssymb, amsthm, enumitem, bbm, calc, graphicx, comment, tikz, t1enc, dsfont, mathrsfs, mathtools, url; the wrapper is the lane's pinned \texttt{amsart} editorial wrapper at 10pt on A4, which restates the theorem-like environments the retained text needs and adds the article's own macro definitions that the retained text needs (\texttt{\textbackslash Emp}, \texttt{\textbackslash MT}, \texttt{\textbackslash eps}), with extra wrapper packages (bm, bbm, dsfont, xspace, cancel, mathtools). The delivered numbering is the unit's own. Assets: no retained span contains a figure environment, a graphics-inclusion command, a PDF-page inclusion, a table or a TikZ picture; no image file is copied into this package, the package declares no asset dependency and no image file is redistributed with it. Licence: version arXiv:2210.08937v1 of this article is distributed under the Creative Commons Attribution 4.0 International licence, as recorded in the arXiv licence metadata for this exact version and shown on the abstract page https://arxiv.org/abs/2210.08937v1. A full rights notice is delivered at licenses/arxiv-2210.08937-CC-BY-4.0.txt. Characters: every retained excerpt is pure ASCII, so the delivered engine, XeLaTeX with its default Latin Modern text font, reports no missing character.
\begin{defn}\label{def:faithful2}
We say that a point $y\in X$ \emph{$(\eps,\delta_1,\delta_2)$-approximately-traces a specification $\xi=(\langle x_j,n_j\rangle)_{j=1}^k$} if
\begin{enumerate}
\item $y\in B_{n_1}(x_1,\eps)$,
\item for each $2\le j \le k$ there is a set $\Lambda_{j}\subseteq \{0,\ldots,n_{j}-1\}$
with $|\Lambda_{j}|\ge (1 - \delta_1) n_{j} $ and an
integer $0\le s_{j-1}\le \delta_2 n_{j}$ such that
\[
T^{(\sum_{i=1}^{j-1}s_i+n_i)}(y)\in B_{\Lambda_j}(x_j,\eps).
\]
\end{enumerate}
\end{defn}

\begin{defn} \label{def:saps}
A Polish dynamical system $(X, T )$ has the \emph{strong approximate product
structure} if the following holds: Given $\eps, \delta_1, \delta_2 > 0$,
there exists a nonegative integer $N=N(\eps,\delta_1,\delta_2)$
such that for any specification $\xi=(\langle x_j,n_j\rangle)_{j=1}^k$ with $n_j \ge N$ for $1 \le j \le k$,
there exists a point $y \in X$ that $(\eps,\delta_1,\delta_2)$-approximately-traces $\xi$.
\end{defn}

\begin{rem}\label{rem:N}
We will also say that a point $y$ $\eps$-approximately-traces a specification $\xi$, if it $(\eps, \eps, \eps)$-approximately-traces $\xi$ as in Definition~\ref{def:faithful2}. We define $N(\eps) = N(\eps, \eps, \eps)$, where $N(\cdot,\cdot,\cdot)$ is as in Definition~\ref{def:saps}. %We will use this abbreviation throughout the proofs.
\end{rem}

\begin{lem}\label{lem:close-segments}
Let $\gamma,\delta,\eps>0$. Assume that
$\bar x=\langle x_1,\ldots,x_k\rangle\in X^k$ and
$\bar y=\langle y_1,\ldots,y_\ell\rangle\in X^\ell$, where $0<k\le \ell\le (1+\delta)k$.
If there exists $j$ with $(1-\gamma)k\le j\le k$ and
subsequences $\bar x'=\langle x_{a_1},\ldots,x_{a_j}\rangle\in X^j$ and
$\bar y'=\langle y_{b_1},\ldots,y_{b_j}\rangle\in X^j$ with $1\le a_1 <\ldots < a_j\le k$
and $1\le b_1 <\ldots < b_j\le \ell$
such that $\rho(x_{a_i},y_{b_i})<\eps$ for $1\le i\le j$, then
\[
D(\Emp(\bar x),\Emp(\bar y))
% \le D\left(\frac{1}{k}\sum_{p=1}^{k}\delta_{x_p},\frac{1}{\ell}\sum_{q=1}^{\ell}\delta_{y_q},\right)
\le 2\gamma+\delta+\eps.
\]
\end{lem}

\begin{lem}\label{lem:protogen}
If $(X,T)$ is a Polish dynamical system with the strong approximate
product structure and $\mu$ is a $T$-invariant measure on $X$, then
for each $\eps>0$ there exists a positive integer $Q$ such that for every $n>0$
%there are $L\in\N$ and $x\in X$ satisfying $nQ\le L< (1+\eps)nQ$ and
there is $x\in X$ satisfying $D(\mu,\Emp(x,nQ))<\eps$.
\end{lem}

\begin{lem} \label{lem:main}
Let $(X, T)$ be a Polish dynamical system with the strong approximate product structure.
\begin{enumerate}
\item \label{it:i} Given a sequence $(\mu_n)_{n \geq 1}$ in $\MT(X)$ such that $D(\mu_n, \mu_{n+1}) \to 0$, there exists a point $y \in X$ and a sequence $L_n \nearrow \infty$ such that
\begin{equation} \label{eqn:lemmain}
D(\Emp(y, j), \mu_{n(j)}) \to 0 \quad \textrm{ as } j \to \infty,
\end{equation}
where $n(j)$ denotes the unique integer satisfying $L_{n(j) - 1} \le j < L_{n(j)}$. In particular, this implies $V(y) = \{ \textrm{ limit points of the sequence } (\mu_k)_{k \geq 1} \}$.
\item \label{it:ii} Given any $y_0 \in X$ and $\eps > 0$, we can define $y$ in \eqref{it:i} so that $y \in B(y_0, \eps)$.
\item \label{it:iii} If $y^{1}$ and $y^{2}$ are points obtained as in \eqref{it:i}  for sequences of measures $( \mu^{1}_n ) _{n \ge 1}$ and $( \mu^{2}_n ) _{n \ge 1}$ such that for some $N\ge 1$ we have that $\mu^{1}_i = \mu^{2}_i$ for $1\le i \le N$,  then $\rho(y^{1}, y^{2}) \le 1 / 2^N$.
\end{enumerate}
\end{lem}

\begin{proof}
Fix $0< \eps < 1$.
We will define sequences of integers $(K_n)_{n\ge 0}$ and points
$(x_n)_{n\ge 0} \subseteq X$ which will be parameters for our construction. For each $n\ge 0$ we use Lemma \ref{lem:protogen}
to pick $x_n \in X$ and $K_n$
so that the following hold:
\begin{gather}
 K_n\ge N(\eps/2^{n+1}), \label{eqn:cond-x}\\
  D(\Emp(x_n,K_n),\mu_n) \le \eps/2^{n+1}.  \label{eqn:cond-i}
\end{gather}
Here, $N(\eps/2^{n+1})$ is defined as in Remark~\ref{rem:N}. %after Definition~\ref{def:saps}.

Next, we will define $(\ell_n)_{n\ge 1}, (L_n)_{n \ge 0}$ and $(y_n)_{n \ge 0} \subseteq X$. We begin with $L_0 = K_1$ and $y_0$ arbitrary point in $X$. For all $n \ge 0$ we will have
\begin{equation} \label{eqn:Ln_bound}
L_{n-1} \ge n K_n.
\end{equation}
We proceed inductively, for $n \ge 1$:
\begin{itemize}
  \item pick $\ell_n$ such that $\ell_n K_n \ge (n + 1) K_{n+1}$ and $\ell_n K_n \ge n L_{n-1}$,
  \item define the specification $\xi_n$ to consist of the pair $(y_{n-1}, L_{n-1})$ followed by the pair $(x_n, K_n)$ repeated $\ell_n$ times,
  \item define $y_n$ to be a point that $\eps/2^{n+1}$-approximately-traces the specification $\xi_n$,
  \item define $L_n$ to be the length of orbit segment required for $y_n$ to trace $\xi_n$. Note that condition (\ref{eqn:Ln_bound}) is satisfied, as $L_n \ge \ell_n K_n \ge n K_{n+1}$.
\end{itemize}
It follows from this definition that
\begin{equation}
\rho(T^i(y_{n-1}), T^i(y_n)) \le \eps / 2^{n+1} \quad \textrm{ for } 0 \le i < L_{n-1}.
\end{equation}
In particular, the sequence $(y_n)$ must converge to a limit $y$, and we have
\begin{equation} \label{eqn:yn_close_y}
\rho(T^i(y_{n-1}), T^i(y)) \le \eps / 2^n \quad \textrm{ for } 0 \le i < L_{n-1}.
\end{equation}

We will now examine the empirical measures $\Emp(y, j)$ for $j>0$. First, consider $\Emp(y, L_n)$ for some integer $n$. By (\ref{eqn:yn_close_y}), we have
\begin{equation} \label{eqn:item1}
D(\Emp(y, L_n), \Emp(y_n, L_n)) \le \eps / 2^n.
\end{equation}
Recall that we have chosen $y_n$ so that it $\eps/2^{n+1}$-approximately-traces $\xi_n$ along the orbit segment of length $L_n$. Together with Lemma \ref{lem:close-segments} this gives $D(\Emp(y_n, L_n), \Emp(\xi_n)) < 4 \eps/2^{n+1} = \eps / 2^{n-1}$.

Define $\bar{\xi_n}$ to be the specification consisting of $\ell_n$ repetitions of pair $(x_n, K_n)$. Thus $\bar{\xi_n}$ only differs from $\xi_n$ by the lack of initial $(y_{n-1}, L_{n-1})$ segment. We defined $\ell_n$ so that $\ell_n K_n \ge n L_{n-1}$, thus an application of Lemma \ref{lem:close-segments} gives
\begin{equation} \label{eqn:item2}
D(\Emp(\xi_n), \Emp(\bar{\xi_n})) \le 1/n.
\end{equation}
Finally, notice that
\begin{equation} \label{eqn:item3}
D(\Emp(\bar{\xi_n}), \mu_n) = D(\Emp(x_n, K_n), \mu_n) \le \eps / 2^{n+1}.
\end{equation}

Combining (\ref{eqn:item1}), (\ref{eqn:item2}) and (\ref{eqn:item3}) gives
\begin{equation} \label{eqn:emp_measure_Ln}
D(\Emp(y, L_n), \mu_n) \le 1/n + 7 \eps/2^{n+1} \to 0 \quad \textrm{ as } n \to \infty.
\end{equation}

We are now prepared to estimate $\Emp(y, j)$ for arbitrary $j>0$. There exists a unique integer $n = n(j)$ such that $L_{n-1} \le j < L_n$. Recall the definition of $y_n$ as a point that $\eps / 2^{n+1}$-approximately-traces $\xi_n$. Let $s_1, \dots s_{\ell_n}$ be the lengths of gaps in this tracing, as in Definition \ref{def:faithful2}. The orbit segment $(y_n, L_n)$ is made up of orbit segments corresponding to the elements of $\xi_n$ (interspersed by gaps of length $s_i$). Thus, the initial subsegment $(y, j)$ fully contains some number $b+1$ of those segments. Formally, $b$ is the unique integer $0 \le b < \ell_n$ such that
\[
B := L_{n-1} + bK_n + \sum_{i < b} s_i \le j < L_{n-1} + (b+1)K_n + \sum_{i \le b} s_i.
\]
We have $j - B \le K_n + s_b \le (1 + \eps / 2^{n+1}) K_n$. We also have
\begin{equation} \label{eqn:B_ineq}
\frac{j - B}{B} \le \frac{(1 + \eps / 2^{n+1}) K_n}{B} \le
\frac{(1 + \eps / 2^{n+1}) K_n}{L_{n-1}} \le \frac{1 + \eps / 2^{n+1}}{n},
\end{equation}
where the last inequality follows from (\ref{eqn:Ln_bound}). Next, note that $y_n$, over the orbit segment of length $B$, $\eps / 2^{n+1}$-approximately-traces the specification
\[
\zeta = \langle
(y_{n-1}, L_{n-1}), \underbrace{(x_n, K_n), \dots (x_n, K_n)}_{b \textrm{ times}}
\rangle.
\]
By the triangle inequality, we obtain
\begin{multline} \label{eqn:temp12}
D(\Emp(y, j), \mu_{n}) \le
D(\Emp(y, j), \Emp(y, B)) + D(\Emp(y, B), \Emp(y_n, B)) + \\
D(\Emp(y_n, B), \Emp(\zeta)) + D(\Emp(\zeta), \mu_{n}).
\end{multline}
Let us examine the terms on the right hand side of $\le$ in \eqref{eqn:temp12}:
\begin{itemize}
  \item $D(\Emp(y, j), \Emp(y, B)) \le (1 + \eps / 2^{n+1}) / n$ by (\ref{eqn:B_ineq}),
  \item $D(\Emp(y, B), \Emp(y_n, B)) \le \eps / 2^{n+1}$ by (\ref{eqn:yn_close_y}),
  \item $D(\Emp(y_n, B), \Emp(\zeta)) \le 4\eps / 2^{n+1}$ by $\eps / 2^{n+1}$-approximate-tracing and Lemma \ref{lem:close-segments},
  \item $D(\Emp(\zeta), \mu_{n}) \le \max \{ D(\Emp(y_{n-1}, L_{n-1}), \mu_{n}), D(\Emp(x_n, K_n), \mu_{n}) \}$.
  Now,
  \begin{itemize}
    \item[$\circ$] $D(\Emp(y_{n-1}, L_{n-1}), \mu_{n-1}) \to 0$ as $n \to \infty$, as shown by (\ref{eqn:emp_measure_Ln}), and by our assumption $D(\mu_n, \mu_{n-1}) \to 0$,
    \item[$\circ$] $D(\Emp(x_n, K_n), \mu_{n}) \le \eps / 2^{n+1} \to 0$ as $n \to \infty$.
  \end{itemize}
\end{itemize}
All of this, applied to the right hand side of $\le$ in \eqref{eqn:temp12}, shows that
\begin{equation}
D(\Emp(y, j), \mu_{n(j)}) \to 0 \quad \textrm{ as } j \to \infty.
\end{equation}
It follows easily that $V(y) \subseteq \{ \textrm{ limit points of the sequence } (\mu_k)_{k \geq 1} \}$. Indeed, if $\Emp(y, j_k) \to \bar{\mu}$ for some subsequence $(j_k)$, then also $\mu_{n(j_k)} \to \bar{\mu}$. Conversely, if $\mu_{n_k} \to \bar{\mu}$ for some $\bar{\mu}$, then $\Emp(y, L_{n_k - 1}) \to \bar{\mu}$.

It remains to address the ``moreover'' part of the lemma statement:
\begin{enumerate}
  \item Regardless of the sequence $(\mu_n)$, by (\ref{eqn:yn_close_y}) we have $\rho(y, y_0) \le \eps / 2$.
  \item Note that construction of $y_1 \dots y_n$ depends only on $\mu_1 \dots \mu_n$. Consider another sequence $(\mu_n')$, such that $\mu_j = \mu_j'$ for $j \le n$. Let $y'$ be the point obtained from this sequence. Then, since $y_n = y_n'$, we deduce from \eqref{eqn:yn_close_y} that
  \[
    \rho(y, y') \le \rho(y, y_n) + \rho(y_n', y') \le \eps /2^{n+1} + \eps /2^{n+1} = \eps / 2^n.
  \qedhere\]
\end{enumerate}
\end{proof}

\end{document}
